% Chapter 1

% Chapter 1

  \chapter{Introducción general}

  \label{Chapter1}
  \label{IntroGeneral}

  El presente capítulo introduce el desarrollo de un sistema multiagente inteligente
  orientado al seguimiento y alerta de activos financieros. Se destaca la
  transformación del sector financiero impulsada por la digitalización y la
  inteligencia artificial (IA), en el que el acceso rápido a la información es clave
  para la toma de decisiones. Las soluciones actuales, como Bloomberg
  \parencite{gratton2024bloomberg,king2016terminal} o FinBERT
  \parencite{araci2019finbert,liu2020finbert,huang2021finbert}, resultan costosas y
  poco accesibles para inversionistas minoristas.

  El trabajo realizado busca reducir la brecha tecnológica y ofrecer una herramienta
  autónoma, modular y de bajo costo que integre análisis predictivo, procesamiento de
  noticias y generación de alertas. Su arquitectura cooperativa de agentes permite
  escalabilidad y personalización del sistema. El impacto esperado es la
  democratización del acceso a la información financiera y la promoción de decisiones
  más informadas y equitativas.

  %----------------------------------------------------------------------------------------   

  \newcommand{\keyword}[1]{\textbf{#1}}
  \newcommand{\tabhead}[1]{\textbf{#1}}
  \newcommand{\code}[1]{\texttt{#1}}
  \newcommand{\file}[1]{\texttt{\bfseries#1}}
  \newcommand{\option}[1]{\texttt{\itshape#1}}
  \newcommand{\grados}{$^{\circ}$}

  %----------------------------------------------------------------------------------------   

  \section{Contexto financiero y tecnológico}

  El ecosistema financiero actual atraviesa un proceso de transformación acelerado
  impulsado por la digitalización, la disponibilidad masiva de datos y los avances en
  IA. En este entorno, el acceso oportuno a la información y la capacidad de
  procesarla en tiempo real se han convertido en factores críticos para la toma de
  decisiones de inversión.

  Al mismo tiempo, los mercados financieros presentan alta volatilidad, asimetrías
  informativas y reacciones tardías ante eventos globales o regionales. Esta dinámica
  evidencia la necesidad de contar con sistemas que analicen datos en tiempo real,
  integren fuentes heterogéneas y ofrezcan señales predictivas para anticipar cambios
  del mercado. Los sistemas multiagentes permiten distribuir tareas entre componentes
  especializados (por ejemplo, agentes de mercado, predicción, noticias o usuario),
  mejorando la velocidad y precisión de las decisiones.

  En particular, la incorporación de modelos de \textit{machine learning} para la
  clasificación de movimientos de precios y de modelos de lenguaje para el análisis
  de sentimiento en noticias financieras representa una tendencia consolidada en la
  literatura reciente \parencite{araci2019finbert,lstm2022financial}. Sin embargo,
  la integración de ambos enfoques en una única plataforma accesible, modular y
  automatizada sigue siendo un desafío abierto, especialmente para usuarios
  minoristas o instituciones financieras de menor escala.

  Como se muestra en la figura~\ref{fig:flujo_informacion_financiera}, el sistema
  transforma datos financieros en alertas personalizadas.

  \begin{figure}[H]
      \centering
      \vspace{-1em}
      \captionsetup{skip=10pt}
      \includegraphics[width=0.9\textwidth]{Figures/Figura1.png}
      \vspace{-1.2em}
      \caption{Flujo de información financiera.}
      \label{fig:flujo_informacion_financiera}
  \end{figure}
  \vspace{-1em}

  %----------------------------------------------------------------------------------------   

  \section{Estado del arte}

  Diversas soluciones comerciales y académicas han abordado el desafío del análisis
  financiero automatizado mediante el uso de tecnologías avanzadas de IA y
  procesamiento de datos. Entre las plataformas comerciales más destacadas se
  encuentra Bloomberg Terminal, una herramienta ampliamente utilizada por
  instituciones financieras que centraliza información de mercado en tiempo real,
  ofrece indicadores económicos, métricas de riesgo y herramientas analíticas para
  la toma de decisiones estratégicas. Sin embargo, su costo elevado y su diseño
  orientado a usuarios expertos la convierten en una solución inaccesible para la
  mayoría de los inversionistas minoristas.

  En el ámbito académico y tecnológico se han desarrollado modelos y arquitecturas
  específicas para el análisis de información financiera. En particular, los modelos
  de lenguaje entrenados en dominios financieros, como FinBERT, aplican técnicas de
  \textit{Natural Language Processing} (NLP) para interpretar el sentimiento presente
  en noticias, reportes corporativos y publicaciones en redes sociales. Este tipo de
  análisis permite cuantificar el tono de la información (positivo, negativo o neutro)        
  y utilizarlo como variable predictiva del comportamiento del mercado.

  Diversos trabajos amplían el análisis de sentimiento en el ámbito financiero. Por
  ejemplo, Amola \parencite{amola2025sentiment} propone un enfoque híbrido que
  integra el análisis de sentimiento de noticias con modelos cuantitativos
  tradicionales. De manera complementaria, estudios más recientes incorporan grandes
  modelos de lenguaje y mecanismos de recuperación de información contextual para
  mejorar la interpretación de textos financieros
  \parencite{llm2023financialsentiment}.

  En paralelo, los modelos de series temporales como LSTM
  \textit{(Long Short-Term Memory)} se utilizan para capturar dependencias temporales
  y estacionales en los datos y proyectar tendencias de precios, volatilidad y volumen        
  de operaciones. Los modelos basados en LSTM han demostrado una alta capacidad para
  capturar dependencias temporales en series financieras
  \parencite{lstm2022financial,comparison2019arima}, e incluso se han combinado con
  análisis de sentimiento proveniente de redes sociales para mejorar la predicción de
  precios \parencite{twitter2024lstm}.

  Más recientemente, los enfoques basados en \textit{ensemble} de clasificadores han
  ganado relevancia en la predicción de la dirección de precios financieros. Métodos
  como Random Forest, Gradient Boosting y XGBoost han demostrado ser competitivos
  frente a modelos de mayor complejidad, especialmente cuando se combinan con un
  proceso riguroso de ingeniería de características técnicas y validación temporal
  \parencite{comparison2019arima}. Este tipo de enfoque, orientado a la clasificación
  binaria del movimiento del mercado (alza o baja), permite obtener señales accionables       
  con mayor interpretabilidad que los modelos de regresión tradicionales.

  No obstante, estas soluciones suelen presentarse de manera fragmentada, donde cada
  modelo o herramienta cumple un propósito aislado dentro del proceso analítico. La
  integración de modelos predictivos, análisis de sentimiento y mecanismos automáticos        
  de alerta en una única plataforma sigue siendo un desafío poco abordado. Además,
  la mayoría de estas implementaciones requieren infraestructura de cómputo de alto
  rendimiento, licencias costosas o conocimientos técnicos avanzados para su
  configuración y mantenimiento.

  Como consecuencia, los inversionistas minoristas y entidades financieras regionales
  enfrentan una desventaja estructural frente a los grandes actores del mercado, al
  no disponer de soluciones tecnológicas asequibles que combinen accesibilidad,
  automatización y capacidad de adaptación. De allí surge la motivación para
  desarrollar un prototipo funcional de sistema inteligente multiagente, capaz de
  integrar estas funciones en una arquitectura modular y escalable, promoviendo así
  un acceso más equitativo al análisis financiero avanzado.

  Como se muestra en la tabla~\ref{tab:comparativa_plataformas}, el prototipo
  propuesto presenta ventajas significativas en costo, accesibilidad e integración
  modular respecto de las soluciones existentes.

  \begin{table}[H]
  \captionsetup[table]{justification=raggedright, singlelinecheck=false}
  \caption{Plataformas existentes vs. el prototipo propuesto.}
  \label{tab:comparativa_plataformas}
  \centering
  {\Large
  \renewcommand{\arraystretch}{1.5}
  \setlength{\tabcolsep}{8pt}
  \resizebox{0.98\textwidth}{!}{%
  \begin{tabular}{
      >{\raggedright\arraybackslash}p{3.8cm}
      >{\raggedright\arraybackslash}p{2.5cm}
      >{\raggedright\arraybackslash}p{2.6cm}
      >{\raggedright\arraybackslash}p{2.8cm}
      >{\raggedright\arraybackslash}p{2.8cm}
      >{\centering\arraybackslash}p{2.5cm}
      >{\raggedright\arraybackslash}p{3.6cm}
  }
  \toprule
  \textbf{Plataforma / Modelo} &
  \textbf{Costo de acceso} &
  \textbf{Accesibilidad} &
  \textbf{Personalización} &
  \textbf{Integración de módulos} &
  \textbf{Enfoque multiagente} &
  \textbf{Tecnologías utilizadas} \\
  \midrule
  \textbf{Bloomberg Terminal} &
  Muy alto (licencia corporativa) &
  Limitada a instituciones financieras &
  Baja (interfaz cerrada) &
  Alta dentro de su ecosistema &
  No &
  API propietaria, bases de datos en tiempo real \\
  \midrule
  \textbf{FinBERT / Modelos NLP financieros} &
  Medio (uso de modelos abiertos) &
  Media (requiere infraestructura técnica) &
  Alta (ajustable por dominio) &
  Media (depende de integración externa) &
  No &
  Python, PyTorch, Transformers \\
  \midrule
  \textbf{Modelos de series temporales (LSTM)} &
  Bajo (software libre) &
  Alta (código abierto) &
  Alta (entrenable con datos propios) &
  Media (requiere ensamblaje manual) &
  No &
  Python, TensorFlow, PyTorch \\
  \midrule
  \textbf{Prototipo propuesto} &
  Bajo / sin costo de licencia &
  Alta (interfaz web accesible) &
  Alta (configurable según usuario) &
  Alta (agentes interconectados) &
  \textbf{Sí} &
  Python, FastAPI, Streamlit, SQLite, ML + NLP \\
  \bottomrule
  \end{tabular}%
  }
  }
  \end{table}

  %----------------------------------------------------------------------------------------   

  \section{Motivación}

  La motivación central de este trabajo radica en la necesidad de reducir la brecha
  tecnológica y de información que separa a las grandes instituciones financieras de
  los pequeños inversionistas. Mientras los actores de gran escala cuentan con
  plataformas integradas, acceso prioritario a datos y capacidad de análisis
  automatizado, los usuarios minoristas suelen depender de información fragmentada,
  tardía o poco contextualizada y muchas veces carecen de conocimientos en el área
  de finanzas. Esta asimetría limita su capacidad para anticipar movimientos del
  mercado y tomar decisiones informadas.

  El desarrollo de un sistema inteligente, autónomo y modular busca revertir esta
  situación mediante la provisión de alertas personalizadas y análisis automatizados,
  que permitan a cada usuario reaccionar de forma oportuna ante eventos financieros
  relevantes. La arquitectura propuesta integra componentes que combinan clasificación        
  de dirección de precios mediante ensemble de modelos, análisis de sentimiento
  proveniente de noticias económicas y visualización de métricas, bajo un enfoque
  accesible, explicable y escalable.

  El enfoque multiagente otorga al sistema flexibilidad y robustez, al permitir que
  cada agente cumpla un rol especializado y opere de manera cooperativa dentro de una
  estructura distribuida. Esta organización facilita la evolución, mantenimiento y
  ampliación del sistema, permitiendo incorporar nuevos módulos sin alterar su
  funcionamiento general.

  La motivación de este trabajo se sustenta en el propósito de democratizar el acceso
  a la inteligencia financiera, brindando a los pequeños inversionistas una herramienta       
  que les permita superar las barreras de conocimiento, tiempo y recursos tecnológicos,       
  y fomentar así un ecosistema financiero más inclusivo, informado y equitativo.

  \section{Objetivos y alcance}

  A continuación, se desarrollan el objetivo general, los objetivos específicos y el
  alcance.

  \subsection{Objetivo general}

  Diseñar e implementar un sistema multiagente inteligente que integre análisis
  predictivo, procesamiento de noticias y generación de alertas personalizadas para
  el seguimiento de activos financieros.

  \subsection{Objetivos específicos}

  Los objetivos específicos son:
  \begin{itemize}
      \item Desarrollar agentes especializados para mercado, predicción, sentimiento,
      recomendación y alertas.
      \item Aplicar un ensemble de clasificadores de \textit{machine learning} para
      predecir la dirección del precio de activos financieros.
      \item Implementar modelos NLP para análisis de sentimiento en noticias financieras,     
      combinando FinBERT, VADER, TextBlob y un lexicón financiero propio.
      \item Diseñar una interfaz gráfica que permita visualizar métricas, predicciones        
      y alertas.
      \item Evaluar el desempeño del sistema mediante métricas de clasificación
      (Accuracy, Precision, Recall, F1-Score, AUC-ROC), latencia y confiabilidad.
  \end{itemize}

  \subsection{Alcance}

  El prototipo se limita al procesamiento de datos públicos y a la generación de
  alertas sobre un conjunto de activos seleccionados. No incluye operaciones reales
  de compra o venta, sino la demostración funcional del modelo multiagente y su
  capacidad de integración. La puesta en producción no forma parte del alcance de
  este trabajo; se propone como línea de trabajo futuro la integración del sistema
  en una aplicación web o móvil de uso institucional o minorista.

  \section{Impacto social y económico}

  La desigualdad informativa en los mercados financieros impacta directamente sobre
  la capacidad de decisión de los pequeños inversionistas. El sistema propuesto
  contribuye a democratizar el acceso a la información y la analítica avanzada,
  ofreciendo:

  \begin{itemize}
      \item Mayor transparencia en el seguimiento de activos financieros.
      \item Reducción de la brecha tecnológica entre grandes y pequeños actores del
      mercado.
      \item Potenciación de la educación financiera a través de herramientas accesibles       
      y explicables.
      \item Fomento de la innovación tecnológica en el ámbito financiero regional.
  \end{itemize}

  Como se observa en la figura~\ref{fig:brecha_informativa}, el sistema reduce la
  brecha informativa entre actores financieros.

  \begin{figure}[H]
      \centering
      \includegraphics[width=0.95\textwidth]{Figures/Figura2.png}
      \caption{Brecha informativa entre grandes inversores y pequeños inversionistas.}        
      \label{fig:brecha_informativa}
  \end{figure}












  

  \section{Organización de la memoria}

  La presente memoria técnica se organiza en cinco capítulos. El capítulo~1 introduce el
  contexto, la motivación, los objetivos y el alcance del sistema desarrollado.
  El capítulo~2 presenta el marco teórico que sustenta las decisiones de diseño,
  incluyendo los fundamentos de los sistemas multiagentes, machine
  learning aplicado a finanzas y el análisis de sentimiento con NLP. El
  capítulo~3 describe en detalle el diseño e implementación del sistema, abarcando
  la arquitectura multiagente, el backend, la base de datos, el sistema de alertas
  y los mecanismos de seguridad. El capítulo~4 presenta los resultados obtenidos
  y la validación del sistema. Finalmente, el capítulo~5 expone las conclusiones
  y las líneas de trabajo futuro.
